\documentclass[11pt]{article}

\usepackage[margin=1.2in]{geometry}
\usepackage{amsmath,amssymb,amsthm}
\usepackage{hyperref}
\usepackage{booktabs}
\usepackage{enumitem}
\usepackage{xcolor}
\usepackage{mdframed}

% ---- theorem-like envs ----
\theoremstyle{definition}
\newtheorem{example}{Example}
\newtheorem{remark}{Remark}
\newtheorem{insight}{Key Insight}

% ---- shaded intuition box ----
\newmdenv[
  backgroundcolor=blue!5,
  linecolor=blue!30,
  roundcorner=4pt,
  skipabove=6pt,
  skipbelow=6pt
]{intubox}

\title{\textbf{Lecture Notes: Community Detection in Networks}\\
\large Modularity and the Louvain Algorithm}
\author{}
\date{}

\begin{document}
\maketitle
\tableofcontents
\newpage

%% ==========================================================
\section{What Is a Community?}

Networks arise everywhere: friendships on a social platform, protein interactions in a cell, routers on the internet. One of the most natural questions to ask about any network is whether its nodes organize into groups---sets of nodes that are tightly connected among themselves but loosely connected to the rest of the network. These groups are called \textbf{communities}.

The intuition is easy to state. A \emph{community} is a subset of nodes with dense internal connections and sparse external connections. In a social network, this might correspond to a circle of friends or colleagues. In a biological network, it might correspond to a functional module of proteins. In both cases, the structure tells us something meaningful about the system that the raw edge list does not make obvious.

Community detection is the task of finding these groups automatically, given only the network. It is useful not just for description---knowing the community structure can guide how you analyze, visualize, or intervene in a network.

%% ==========================================================
\section{Why Simply Counting Internal Edges Is Not Enough}

A first attempt at formalizing ``dense internal connections'' might be: find a partition of nodes that maximizes the number of edges falling \emph{within} communities. The problem is that this naive criterion is trivially satisfied by putting all nodes into a single community---every edge becomes an internal edge.

More subtly: even in a dense network with no community structure at all, a partition into two groups will have many internal edges just by chance. The number of internal edges by itself says nothing about whether the partition is meaningful.

What we really want to know is not how many internal edges there are, but how many \emph{more} than we would expect if the edges were placed at random. This calls for a \textbf{null model}---a reference against which to measure the observed network.

%% ==========================================================
\section{The Null Model: The Configuration Model}

The standard null model for community detection is the \textbf{configuration model}. The idea is simple. Imagine taking the original network and cutting every edge in the middle, leaving two ``stubs'' dangling at each endpoint. Node $i$ now has $k_i$ stubs, where $k_i$ is its degree (or weighted degree, if edges have weights). Then randomly re-pair all the stubs to form new edges. The result is a random network in which every node has the same degree as in the original, but all higher-order structure---including communities---has been destroyed.

\begin{intubox}
Think of it as a randomized rewiring. You preserve the popularity of each node (its degree), but you shuffle who connects to whom. Any community-like pattern that survives this shuffling is trivially explained by degrees alone and should not count as real structure.
\end{intubox}

Under this model, what is the expected number of edges between nodes $i$ and $j$? Node $i$ has $k_i$ stubs. The total number of stubs in the entire network is $2m$, where $m$ is the total edge weight (each edge contributes two stubs). The probability that any one stub of node $i$ pairs with any one particular stub of node $j$ is $1/(2m)$. There are $k_j$ stubs at node $j$, so the probability that a specific stub of $i$ connects to \emph{any} stub of $j$ is $k_j / (2m)$. Multiplying by the $k_i$ stubs at node $i$ gives the expected edge weight between $i$ and $j$:
\[
  \mathbb{E}[A_{ij}] = \frac{k_i k_j}{2m}.
\]
This is the key formula. Notice that high-degree nodes are expected to share many edges even in the randomized network---the null model already accounts for their popularity. Community structure is only the part that goes \emph{beyond} this expectation.

%% ==========================================================
\section{Modularity $Q$}

We now have everything we need to define \textbf{modularity}. Given a partition of nodes into communities, modularity $Q$ measures how much the intra-community edge weight exceeds what the configuration model predicts:

\[
  Q = \frac{1}{2m} \sum_{i,j} \left[ A_{ij} - \frac{k_i k_j}{2m} \right] \delta(c_i, c_j).
\]

\paragraph{Notation.}
\begin{itemize}[leftmargin=2em]
  \item $A_{ij}$: the edge weight between nodes $i$ and $j$ (zero if no edge).
  \item $k_i = \sum_j A_{ij}$: the weighted degree of node $i$.
  \item $m = \frac{1}{2} \sum_{i,j} A_{ij}$: the total edge weight of the network.
  \item $c_i$: the community that node $i$ belongs to.
  \item $\delta(c_i, c_j)$: the Kronecker delta, equal to 1 if $c_i = c_j$ (same community) and 0 otherwise.
\end{itemize}

\paragraph{Interpretation.} The sum runs over all pairs of nodes, but the $\delta$ function restricts it to pairs in the same community. For each such pair, we compute the ``surprise'': how much the actual edge weight $A_{ij}$ exceeds the random expectation $k_i k_j / (2m)$. The factor $1/(2m)$ normalizes the result to the range $[-1, 1]$.

A positive $Q$ means that communities have more internal edge weight than a random graph with the same degree sequence would. $Q = 0$ means the partition is no better than random. $Q < 0$ means communities have \emph{fewer} internal edges than expected, which typically does not arise in practice.

\begin{insight}
  Modularity is not a definition of community---it is one particular formalization. Different null models or different objectives can lead to different notions of community. $Q$ is the most widely used because it is simple, interpretable, and admits efficient optimization algorithms.
\end{insight}

\paragraph{Practical threshold.} In practice, values of $Q$ above roughly 0.3 are taken as evidence of meaningful community structure. Values near 0 suggest the partition is no better than random, and negative values are unusual.

%% ==========================================================
\section{Why $Q$ Is Useful but Imperfect}

Modularity is a practical and widely used tool, but it has a known limitation worth flagging: the \textbf{resolution limit}.

When optimizing $Q$, the algorithm tends to merge small communities into larger ones if the small communities are not large enough relative to the total network size. Concretely, two communities that are genuinely well-separated may be lumped together if each is small compared to $m$. This means that $Q$ optimization can miss structure at fine scales in large networks.

This is not a reason to abandon $Q$---it remains extremely useful in practice---but it is a reason to interpret results critically, especially when communities of very different sizes are expected.

%% ==========================================================
\section{The Louvain Algorithm}

Finding the partition that exactly maximizes $Q$ is computationally intractable: the number of possible partitions of $N$ nodes grows faster than exponentially, making exhaustive search impossible for any real network. We need a fast heuristic.

The \textbf{Louvain algorithm} (Blondel et al., 2008) is the most widely used method for approximately maximizing $Q$. It is fast, scales to networks with millions of nodes, and produces high-quality partitions in practice. The algorithm alternates between two phases, repeated until convergence.

\subsection{Phase 1: Local Moves}

Start by assigning each node to its own community---$N$ communities of size one. Then, for each node $i$, consider moving it into the community of each of its neighbors. For each candidate move, compute the change in $Q$ that would result. Move node $i$ to whichever neighboring community gives the largest \emph{positive} gain in $Q$. If no move increases $Q$, leave node $i$ where it is.

Sweep through all nodes in this way, repeatedly, until a full pass through all nodes produces no further improvement. At this point, Phase~1 has reached a local optimum.

\begin{intubox}
The key insight is that the modularity gain from moving a single node can be computed efficiently using only local information: the degrees of the node and its neighbors, and the total degree of the destination community. This makes each move cheap, and the full sweep fast even for large networks.
\end{intubox}

\subsection{Phase 2: Aggregation}

Once Phase~1 converges, \textbf{collapse each community into a single super-node}. Build a new coarsened network where:
\begin{itemize}[leftmargin=2em]
  \item Each super-node represents one community from Phase~1.
  \item The edge weight between two super-nodes is the sum of all edge weights between the corresponding communities in the original network.
  \item Edges within a community become \emph{self-loops} on the corresponding super-node, preserving their contribution to the degree.
\end{itemize}
The result is a smaller network with the same total edge weight $m$, but with community-level granularity.

\subsection{Iteration}

Run Phase~1 again on the coarsened network, then Phase~2 again, and so on. Each round further merges communities, producing a coarser partition. The algorithm stops when a full pass of Phase~1 on the current network produces no improvement in $Q$.

\paragraph{Hierarchy.} A natural byproduct of this iterative aggregation is a \textbf{hierarchical decomposition}: the communities found at each level of aggregation form a nested tree. The finest level has many small communities; each subsequent level merges them into progressively larger ones. This hierarchy can be valuable in its own right, revealing structure at multiple scales.

\paragraph{Complexity.} For sparse networks, the algorithm runs in time effectively linear in the number of edges. Each Phase~1 sweep costs time proportional to the number of edges, and the number of sweeps until convergence is small in practice---typically a handful of passes even on large networks.

\paragraph{Convergence and quality.} Louvain is not guaranteed to find the global maximum of $Q$---it is a greedy local search. In practice, it consistently finds high-quality partitions and is fast enough to run multiple times with different random orderings to check robustness.

%% ==========================================================
\section{Summary}

The chain of ideas in these notes has a simple logic. Communities are an intuitive concept---dense groups with sparse external connections---but making this precise requires care. Simply counting internal edges is misleading without a baseline. The configuration model provides that baseline, giving us the expected edge weight between any pair of nodes under a degree-preserving null. Modularity $Q$ measures how much a given partition beats this baseline. And the Louvain algorithm is a fast, practical method for finding partitions with high $Q$, iterating between local greedy moves and network aggregation until no further improvement is possible.

Together, these tools form a practical and principled pipeline for discovering structure in networks.

\bigskip
\noindent\textit{Reference: V.~D.~Blondel, J.-L.~Guillaume, R.~Lambiotte, E.~Lefebvre, ``Fast unfolding of communities in large networks,'' Journal of Statistical Mechanics: Theory and Experiment, 2008(10):P10008.}

\end{document}
